\appendix
\section{The exact instructions given to the proposer}
\label{app:prompts}

Every claim in this paper about what the model was and was not told rests on the
text below. These are the task messages verbatim, as sent to the proposer at
every generation, reproduced in full so the reader can check them rather than
take our word for it. Nothing else describes the task to the model: the only
other inputs are the seed program, the formal gate schema included below, and the
numerical feedback from each evaluation.

\subsection{Anonymised motif task (Sections 5--6)}

The claim is that this text never names tic-tac-toe and never mentions boards,
rows, columns, diagonals, corners, centres or winning lines. It does disclose the
qubit connectivity, because that constrains which circuits are legal, and it does
name the readout groups, which is the intentional partial-symmetry disclosure
discussed in Section~\ref{sec:anon}.

\begin{lstlisting}[style=prompt]
You are an expert in quantum machine learning and variational quantum circuits.

Goal:
Evolve the ANSATZ_SPEC for a 9-qubit, 3-class classifier.  The fixed code will
train each candidate with Adam in a simulated PennyLane circuit and report
accuracy, loss, and generalization metrics.

Fixed architecture:
- 9 qubits, indexed 0..8.
- Feature map: RX(2*pi/3 * x_i) on qubit i, for an input vector x in {-1,0,+1}^9.
- Data re-uploading l=3.
- The candidate ansatz block is applied p=2 times after each re-uploading.
- The block receives independent parameter copies for each upload/repetition.
- The three class scores are fixed linear readouts of the final per-qubit
  Pauli-Z expectations.

Only edit the EVOLVE-BLOCK, especially ANSATZ_SPEC.  Do not change the fixed
training loop, data, readout, l, p, or feature map.

Formal ANSATZ_SPEC schema:
- Single-qubit parametrized gates:
  {"gate": "RX"|"RY"|"RZ", "wire": int 0..8, "param": "name"}
- Fixed two-qubit gates:
  {"gate": "CNOT"|"CZ", "wires": [first, second]}
- Parametrized controlled rotations:
  {"gate": "CRX"|"CRY"|"CRZ", "wires": [control, target], "param": "name"}
- Parametrized three-qubit interactions (any three distinct qubits 0..8):
  {"gate": "ZZZ", "wires": [a, b, c], "param": "name"}   # collective exp(-i*theta/2 * Z@Z@Z)
  {"gate": "CCRZ", "wires": [control_1, control_2, target], "param": "name"}   # doubly-controlled RZ

Connectivity constraint:
Two-qubit gates are allowed only on these hardware-native qubit pairs:
(0,2), (0,3), (0,7), (0,8), (1,3), (1,8),
(2,4), (2,6), (3,5), (4,8), (5,7), (6,7).
Three-qubit gates have no connectivity restriction: any 3 distinct qubits.

Parameter sharing:
Reusing the same param string shares that parameter within one ansatz block.
Use sharing deliberately for parameter efficiency.

Starting point:
The seed is an EfficientSU2-style block: independent RY rotations on every
qubit, independent pre-entanglement RZ rotations, CZ gates on each
hardware-native pair, and independent post-entanglement RZ rotations.

Candidate quality:
Improve validation/test accuracy, reduce the train-test gap, reduce L2 loss,
use parameters efficiently, and reach high accuracy in fewer Adam steps.

Invalid candidates will be rejected if they use unsupported gates, bad wires,
non-native two-qubit operations, non-finite metrics, or too many parameters.
\end{lstlisting}

\subsection{Network task (Section~\ref{sec:sym}, task A)}

This text never mentions symmetry, permutation, invariance, equivariance, orbits,
groups, or graphs, and never states that relabelling the qubits leaves the label
unchanged. It does state the premise the model reasons from, in the fourth line
of the fixed architecture: each of the 28 features is a coupling on a qubit pair,
drawn from a pair table. The seed program's docstring repeats it as
\texttt{FEATURE\_PAIRS[k]}, and ShinkaEvolve shows the seed in full. From there
$28 = \binom{8}{2}$ is arithmetic and the graph reading follows. We originally
described this text as withholding the correspondence, which was wrong; the
correction and the run that establishes it are in Section~\ref{sec:zcsym}.

\begin{lstlisting}[style=prompt]
You are an expert in variational quantum circuit design.

Goal:
Evolve the ANSATZ_SPEC for an 8-qubit classifier of binary-feature records.
The fixed code will train each candidate with Adam in a simulated PennyLane
circuit.

Fixed architecture:
- 8 qubits.
- Each record has 28 binary features; the fixed feature map encodes feature k
  as a two-qubit phase coupling (IsingZZ) on a fixed qubit pair given by the
  dataset's pair table.
- Data re-uploading l=3; the candidate ansatz block is applied p=2 times after
  each upload, with independent parameter copies per block.
- Readout: the mean of all single-qubit Z expectations, passed through a
  trainable gain and bias, regressed to the +1/-1 label.

Only edit the EVOLVE-BLOCK, especially ANSATZ_SPEC. Do not change the fixed
training loop, dataset handling, measurements, l, p, or feature map.

Formal ANSATZ_SPEC schema:
- Single-qubit parametrized gates:
  {"gate": "RX"|"RY"|"RZ", "wire": int 0..7, "param": "name"}
- Fixed two-qubit gates:
  {"gate": "CNOT"|"CZ", "wires": [first, second]}
- Parametrized controlled rotations:
  {"gate": "CRX"|"CRY"|"CRZ", "wires": [control, target], "param": "name"}

Connectivity: two-qubit gates may act on ANY pair of distinct qubits.

Parameter sharing:
Reusing the same param string shares that parameter within one ansatz block.

Starting point:
The seed block applies independent RY and RZ rotations on every qubit, a line
of CZ gates, and a final RZ layer.

Candidate quality:
Improve validation/test accuracy, reduce the train-test gap, reduce L2 loss,
use parameters efficiently, and reach high accuracy in fewer Adam steps.

Invalid candidates are rejected for unsupported gates, bad wires, non-finite
metrics, or too many parameters.
\end{lstlisting}

\subsection{Quantum-state task (Section~\ref{sec:sym}, task B)}

The claim is that this text never mentions spin, SU(2), rotation invariance,
isotropic exchange, the Heisenberg model, singlets or dimerisation. The inputs
are described only as ``a precomputed 8-qubit state vector''. The one
non-neutral statement is that the score rewards parameter economy, which is a
property of the fitness function and not a hint about which structure achieves
it.

\begin{lstlisting}[style=prompt]
You are an expert in variational quantum circuit design.

Goal:
Evolve the ANSATZ_SPEC for an 8-qubit classifier of quantum states.
The fixed code will train each candidate with Adam in a simulated PennyLane
circuit.

Fixed architecture:
- 8 qubits.
- Each input is a precomputed 8-qubit state vector from the dataset, prepared
  with StatePrep.
- The candidate ansatz block is applied 6 times in sequence, with independent
  parameter copies per repetition.
- Readout: the mean of a fixed set of two-qubit Z-Z correlators given by the
  dataset, passed through a trainable gain and bias, regressed to the +1/-1
  label.

Only edit the EVOLVE-BLOCK, especially ANSATZ_SPEC. Do not change the fixed
training loop, dataset handling, measurements, or block count.

Formal ANSATZ_SPEC schema:
- Single-qubit parametrized gates:
  {"gate": "RX"|"RY"|"RZ", "wire": int 0..7, "param": "name"}
- Fixed two-qubit gates:
  {"gate": "CNOT"|"CZ", "wires": [first, second]}
- Parametrized two-qubit gates:
  {"gate": "CRX"|"CRY"|"CRZ", "wires": [control, target], "param": "name"}
  {"gate": "XX"|"YY"|"ZZ", "wires": [a, b], "param": "name"}   # Ising-type rotations

Connectivity: two-qubit gates may act on ANY pair of distinct qubits.

Parameter sharing:
Reusing the same param string shares that parameter within one ansatz block.

Starting point:
The seed block applies independent RY and RZ rotations on every qubit, a line
of CZ gates, and a final RZ layer.

Candidate quality:
The scoring on this task strongly rewards PARAMETER ECONOMY and fast
convergence: match or beat the baseline accuracy with as few trainable
parameters as possible. Also improve validation/test accuracy, reduce the
train-test gap, and reduce L2 loss.

Invalid candidates are rejected for unsupported gates, bad wires, non-finite
metrics, or too many parameters.
\end{lstlisting}

\subsection{Zero-context version of the network task (Section~\ref{sec:zcsym})}

This is the same task with the feature-to-pair correspondence removed. The
message states the qubit count, the gate schema, the sharing mechanism and what
the score rewards, and nothing else. The number 28 does not appear, so the
arithmetic that led to the graph reading is unavailable.

\begin{lstlisting}[style=prompt]
You are optimizing a parametrized quantum circuit.

Goal:
Evolve the ANSATZ_SPEC below. Fixed code trains each candidate and returns
numeric feedback. You are NOT told what the data is, what the inputs represent,
what the labels mean, or how the inputs are encoded into the circuit, and you do
not need to know any of it. Your only guide is the feedback you receive.

Fixed architecture:
- 8 qubits, indexed 0..7.
- The input encoding, the readout, the training loop and the metrics are fixed
  and are implemented outside this file.
- The candidate ansatz block is applied a fixed number of times, with
  independent parameter copies per repetition.

Only edit the EVOLVE-BLOCK, that is ANSATZ_SPEC. Nothing else is editable.

Formal ANSATZ_SPEC schema:
- Single-qubit parametrized gates:
  {"gate": "RX"|"RY"|"RZ", "wire": int 0..7, "param": "name"}
- Fixed two-qubit gates:
  {"gate": "CNOT"|"CZ", "wires": [first, second]}
- Parametrized controlled rotations:
  {"gate": "CRX"|"CRY"|"CRZ", "wires": [control, target], "param": "name"}

Two-qubit gates may act on ANY pair of distinct qubits.

Parameter sharing:
Reusing the same param string shares that parameter within one ansatz block.
Use sharing deliberately: it is the only way to reduce the parameter count
without removing gates.

Candidate quality:
Improve validation accuracy, reduce the train-test gap, reduce
L2 loss, use parameters efficiently, and converge in fewer steps.

Invalid candidates are rejected for unsupported gates, bad wires, non-finite
metrics, or too many parameters.
\end{lstlisting}

\noindent Because the seed file is shown in full, it was rewritten too. Every
line describing the data, the encoding, the readout, the training loop or the
metrics moved into an imported module. What remains above the editable region is:

\begin{lstlisting}[style=prompt]
"""Seed program. Only ANSATZ_SPEC inside the EVOLVE-BLOCK is evolved.

Everything else about the task, that is how inputs are encoded, how the circuit
is measured, how training works and how metrics are computed, is fixed and lives
in a module that is not reproduced here. No information about the data is
available in this file.
"""

from _backend import run_experiment as _run

N_QUBITS = 8
ALLOWED_SINGLE_QUBIT_GATES = {"RX", "RY", "RZ"}
ALLOWED_TWO_QUBIT_GATES = {"CNOT", "CZ"}
ALLOWED_PARAM_TWO_QUBIT_GATES = {"CRX", "CRY", "CRZ"}
\end{lstlisting}

\noindent For comparison, the hinted condition's seed docstring reads: ``Records
are 28-dimensional binary feature vectors from an external dataset. Each feature
$k$ is encoded by the fixed feature map as a two-qubit phase coupling on the
qubit pair \texttt{FEATURE\_PAIRS[k]}.'' That is the sentence whose removal
costs the result.

\subsection{For contrast: what the leaked run was shown}

The following sits at the top of the seed program of the original run
(Section~\ref{sec:leaky}), above the editable region, and is therefore included
in the proposer's context in full at every generation. It is reproduced to make
the difference between the two conditions concrete.

\begin{lstlisting}[style=prompt]
# Paper indexing:
#   0 -- 1 -- 2
#   |    |    |
#   7 -- 8 -- 3
#   |    |    |
#   6 -- 5 -- 4
PAPER_GRID = ((0, 1, 2), (7, 8, 3), (6, 5, 4))
CORNERS = (0, 2, 4, 6)
EDGES = (1, 3, 5, 7)
CENTER = 8
GRID_EDGES = (
    (0, 1), (1, 2), (2, 3), (3, 4),
    (4, 5), (5, 6), (6, 7), (7, 0),
    (1, 8), (3, 8), (5, 8), (7, 8),
)
GRID_EDGE_SET = {tuple(sorted(edge)) for edge in GRID_EDGES}
WIN_LINES = (
    (0, 1, 2), (7, 8, 3), (6, 5, 4),
    (0, 7, 6), (1, 8, 5), (2, 3, 4),
    (0, 8, 4), (2, 8, 6),
)
\end{lstlisting}

\noindent The eight winning lines are listed outright, together with the board
geometry, the corner and centre roles, and an ASCII diagram of the grid. No
search was required to recover any of it.
